\section{Presentation Day}

This section explains what is expected on the final presentation day. Each group will deliver one presentation that summarises the project problem, the proposed solution, and the main lessons learned during the semester. Whenever possible, share speaking responsibility across the team so that the presentation reflects collective ownership of the work.

The standard format is:
\begin{itemize}
    \item \textbf{10 minutes} for the main presentation; and
    \item \textbf{5 minutes} for questions and discussion.
\end{itemize}

Electronic slides are the default presentation format. Live demonstrations may be used if they are reliable, but pre-prepared slides, images, or recorded demonstrations are usually safer because they reduce the risk of technical failure.

The presentation will normally be evaluated on the clarity and professionalism of the communication rather than on presentation style alone. A strong presentation should:
\begin{itemize}
    \item define the problem clearly;
    \item show how customer needs and specifications informed the concept;
    \item explain the technical and design rationale behind the proposed solution; and
    \item reflect briefly on the strengths, limitations, and next steps of the project.
\end{itemize}

I plan to make the presentation infront of industry  guests. Therefore, 
it is important to have interesting product concepts. Also, it is important to have a clear and concise presentation that effectively communicates the problem, solution, and learning outcomes. The presentation should be well-structured, visually appealing, and engaging for the audience. It should also demonstrate a deep understanding of the project and the applied physics concepts involved.